% Ontologizer architecture figures.
%
% Requires only what paper.tex already loads:
%   tikz + \usetikzlibrary{arrows.meta, positioning}, amsmath,
%   caption, subcaption.
%
% Values annotated on the diagrams are the live SONAR configuration
% (sonar.py): d = 1024, e_dec = 2048, k = 32, h = 32, l = 5, n = 2
% (bilinear), forward = "resid", deepsup = True, resid_norm = True,
% resid_const = True, encoded = False, scaled = False.

% =====================================================================
% (1) the layer stack
% =====================================================================
\begin{figure}[tp]
\centering
\resizebox{\linewidth}{!}{\begin{tikzpicture}[
    font=\footnotesize,
    >={Stealth[round,length=1.8mm,width=1.6mm]},
    blk/.style={draw, rounded corners=2pt, align=center,
                minimum height=8mm, inner xsep=4pt},
    hook/.style={draw, rounded corners=2pt, densely dashed, align=center,
                 inner sep=3pt, font=\scriptsize, black!55},
    dec/.style={draw, rounded corners=2pt, fill=black!8,
                minimum height=6.5mm, inner xsep=3pt},
    cond/.style={draw, rounded corners=2pt, fill=black!4, align=center,
                 inner sep=2.5pt, font=\scriptsize},
    add/.style={draw, circle, minimum size=4.5mm, inner sep=0pt},
    lab/.style={inner sep=1.5pt, font=\scriptsize},
    sig/.style={->, semithick},
    aux/.style={->, densely dashed, semithick, black!45},
  ]

  %% ---- layer row --------------------------------------------------
  \node[lab] (x)  at (-1.5,0) {$\mathbf{x}\in\mathbb{R}^{d}$};
  \node[blk] (D0) at ( 0.3,0) {$\mathtt{DictEnc}_1$};
  \node[blk] (D1) at ( 4.3,0) {$\mathtt{DictEnc}_2$};
  \node      (Dd) at ( 7.8,0) {$\cdots$};
  \node[blk] (DL) at ( 9.7,0) {$\mathtt{DictEnc}_{l}$};

  \draw[sig] (x) -- (D0);

  \node[hook] (noise) at (-1.5,2.6) {$+\,\varepsilon$};
  \draw[aux] (noise) -- (x);

  %% ---- residual accumulator (top rail, y = 2.6) --------------------
  \node[add] (A0) at ( 2.3,2.6) {$+$};
  \node[add] (A1) at ( 6.1,2.6) {$+$};
  \node[add] (AL) at (11.9,2.6) {$+$};

  \node[lab] (R0) at (0.6,2.6) {$R_0\!=\!0$};
  \draw[sig]      (R0) -- (A0);
  \draw[sig]      (A0) -- (A1);
  \draw[sig]      (A1) -- (AL);
  \draw[semithick] (AL) -- (12.35,2.6);

  %\node[lab] at (8.9,3.05) {$R\in\mathbb{R}^{e}$: residual accumulator};

  %% ---- per-layer contributions F_i (enter from below) --------------
  \draw[sig] (D0.north) -- ++(0,1.1) -| (A0.south);
  \draw[sig] (D1.north) -- ++(0,1.1) -| (A1.south);
  \draw[sig] (DL.north) -- ++(0,1.1) -| (AL.south);
  \node[lab] at (0.62,1.05) {$F_1$};
  \node[lab] at (4.62,1.05) {$F_2$};
  \node[lab] at (10.05,1.05) {$F_{l}$};

  %% ---- shared decoder, one decode per prefix -----------------------
  \node[dec] (Y1) at ( 2.75,-1.9) {$\mathrm{dec}$};
  \node[dec] (Y2) at ( 6.55,-1.9) {$\mathrm{dec}$};
  \node[dec] (YL) at (12.35,-1.9) {$\mathrm{dec}$};

  \draw[sig] ( 2.75,2.6) -- (Y1.north);
  \draw[sig] ( 6.55,2.6) -- (Y2.north);
  \draw[sig] (12.35,2.6) -- (YL.north);
  \fill ( 2.75,2.6) circle (1.1pt);
  \fill ( 6.55,2.6) circle (1.1pt);

  \node[lab,anchor=west] at ( 2.85,1.0) {$R_1$};
  \node[lab,anchor=west] at ( 6.65,1.0) {$R_2$};
  \node[lab,anchor=west] at (12.45,1.0) {$R_l$};

  %% ---- deep supervision --------------------------------------------
  \node[lab] (L1) at ( 0.9,-1.9) {$\mathrm{MSE}(\mathbf{x},\hat{\mathbf{x}}_1)$};
  \node[lab] (L2) at ( 4.9,-1.9) {$\mathrm{MSE}(\mathbf{x},\hat{\mathbf{x}}_2)$};
  \node[lab] (LL) at (10.7,-1.9) {$\mathrm{MSE}(\mathbf{x},\hat{\mathbf{x}}_l)$};
  \draw[aux] (Y1.west) -- (L1);
  \draw[aux] (Y2.west) -- (L2);
  \draw[aux] (YL.west) -- (LL);

  %% ---- residual forwarding ------------------------------------------
  \node[cond] (C1) at (2.75,-3.9)
  {$\mathrm{sg}\,(\mathbf{x}-\hat{\mathbf{x}}_1)$\\ normalise, $\oplus\,1$};
  \node[cond] (C2) at (6.55,-3.9)
  {$\mathrm{sg}\,(\mathbf{x}-\hat{\mathbf{x}}_2)$\\ normalise, $\oplus\,1$};
  \draw[sig] (Y1) -- node[right,lab]{$\hat{\mathbf{x}}_1$} (C1);
  \draw[sig] (Y2) -- node[right,lab]{$\hat{\mathbf{x}}_2$} (C2);

  \draw[sig] (C1.east) -- node[left,lab,pos=0.72]{$\tilde{\mathbf{y}}_1$} (D1.south);
  \draw[sig] (C2.east) -- node[left,lab,pos=0.72]{$\tilde{\mathbf{y}}_2$} (7.8,-0.5);

  %% ---- model output --------------------------------------------------
  \node[lab] (out) at (12.35,-3.9) {$\hat{\mathbf{x}}$};
  \draw[sig,thick] (YL) -- node[right,lab]{$\hat{\mathbf{x}}_l$} (out);

  %% ---- x broadcast to every conditioning step -------------------------
  \draw[aux] (x.south) -- (-1.5,-5.3) -- (2.75,-5.3) -- (C1.south);
  \draw[aux] (2.75,-5.3) -- (6.55,-5.3) -- (C2.south);
  \fill[black!45] (2.75,-5.3) circle (1.1pt);

\end{tikzpicture}

}
\caption[\texttt{Ontologizer} architecture]{\textbf{\texttt{Ontologizer} architecture.}
	An \texttt{Ontologizer} consists of $l$ stacked \texttt{DictEnc} layers.
	Each returns a feature vector $F_i\in\mathbb{R}^{e}$,
	which is written to a residual accumulator $R$.
	For each layer $i$, $R_i = R_{i-1} + F_i$; a stop-gradient variant
	($\mathrm{sg}$ on $R_{i-1}$, available but off in the live
	configuration) would prevent gradient communication between layers.
	$\mathrm{dec}$ is a linear decoder with shape $\mathbb{R}^{d \times e}$.
	It uses deep supervision to obtain a coarse-to-fine representation where
	the component encoded by previous layers is inaccessible to later layers.
	At each layer $i$, $\hat{\mathbf{x}}_i = \mathrm{dec}(R_i)$.
	The same decoder is shared by all layers.
	The component left unexplained after layer $i$ is given by
	$\mathbf{y}_i = \mathbf{x} - \hat{\mathbf{x}}_i$.
	The input to layer $i+1$ is given by
	$\tilde{\mathbf{y}}_i
	  = \mathbf{y}_i / \mathrm{L2}(\mathbf{y}_i) \oplus [1]$
	where $\mathrm{L2}$ is the unit norm. This prevents exploding gradients from repeated
	application of bilinear layers, as each scales its input by $\lVert\cdot\rVert^{2}$.
	A constant coordinate $[1]$ is appended, giving a dimension of $d+1$.
	This prevents the sign information of $\mathbf{y}_i$ from being lost
	by the bilinear transformation.
	Effectively, it adds a bias to an otherwise unbiased linear form
	while preserving eigendecomposition.

	$\mathrm{MSE}(\mathbf{x},\mathrm{dec}(R_i))$ is computed separately for each layer.

	The input-noise hook drawn on $\mathbf{x}$ is available but disabled
	in the live configuration.
}
\label{fig:onto-stack}
\end{figure}

% =====================================================================
% (2) inside one DictEnc, and the per-head geometry
% Split into two floats: the combined float (with both long captions)
% overflowed the page by ~74pt, and its wrapper label was never
% referenced.
% =====================================================================
\begin{figure}[tp]
\centering
\resizebox{\linewidth}{!}{\begin{tikzpicture}[
    font=\footnotesize,
    >={Stealth[round,length=1.8mm,width=1.6mm]},
    blk/.style={draw, rounded corners=2pt, align=center,
                minimum height=8mm, inner xsep=4pt},
    hook/.style={draw, rounded corners=2pt, densely dashed, align=center,
                 inner sep=3pt, font=\scriptsize, black!55},
    lab/.style={inner sep=1.5pt, font=\scriptsize, align=center},
    sig/.style={->, semithick, align=center},
    aux/.style={->, densely dashed, semithick, black!45},
  ]

  %% ---- main chain ----------------------------------------------------
  \node[lab] (E)  at (-0.6,0) {$\tilde{\mathbf{y}}$};
  \node[blk] (CL) at ( 1.2,0) {$\bilinearblock^{d \to h \times k}$};
  \node[blk] (SM) at ( 4.9,0) {$\mathrm{softmax}_k(\cdot/t)$};
  \node[blk] (DB) at ( 9.4,0) {$\dictblock^{h \times k \to h \times e}$};
  \node[blk] (SU) at (12.3,0) {$\sum^{h}_{i=1}$};
  \node[lab] (F)  at (13.6,0) {$F'$\\ pooled\\ features};

  \draw[sig] (E)  -- (CL);
  \draw[sig] (CL) -- node[below,lab]{$K$\\ logits}   (SM);
  \draw[sig] (SM) -- node[below,lab]{$P$\\ classifications}   (DB);
  \draw[sig] (DB) -- node[below,lab]{$F$\\ features} (SU);
  \draw[sig] (SU) -- (F);

  %% ---- training-time hooks and the intervention interface (above) ----
  \node[hook] (nk) at (3.0,1.5) {$+\,\varepsilon_K$,\ winner dropout};
  \node[hook] (iv) at (7.0,2.3) {\texttt{DictIntervention}\\
                                 set / add / sub / ablate / scale};
  \node[hook] (nf) at (11.4,1.5) {$+\,\varepsilon_F$};
  \draw[aux] (nk) -- (3.0,0.20);
  \draw[aux] (iv) -- (7.0,0.20);
  \draw[aux] (nf) -- (11.4,0.20);

  %% ---- where each regulariser attaches (below) -----------------------
  \node[hook] (lp) at (4.9,-1.8)
    {$H$\quad $\mathrm{KL}_m$\quad $\mathrm{cossim}_b$};
  \node[hook] (lf) at (9.4,-1.8)
    {$\mathrm{L1}$\quad $\sigma$};
  \draw[aux] (SM.south) -- (lp);
  \draw[aux] (DB.south) -- (lf);

\end{tikzpicture}

}
\caption[\texttt{DictEnc} architecture]{\textbf{\texttt{DictEnc} architecture.}
	Dashed lines above indicate perturbations to the data,
	including noising and causal interventions.
	Dashed lines below indicate statistics used as tuneable auxiliary loss functions.
	Each \texttt{DictEnc} consists of a multihead bilinear classifier
	($\bilinearblock$) with weights
	$[\mathsf{W}, \mathsf{V}] \in \mathbb{R}^{2 \times h \times k \times d}$
	($d+1$ after the first layer, per the appended constant coordinate)
	and a multihead linear dictionary ($\dictblock$)
	with weights $\mathsf{D} \in \mathbb{R}^{h \times k \times e}$.
	It accepts a vector $\mathbf{x} \in \mathbb{R}^{d}$,
	which may be a residual from a previous layer.
	A $\bilinearblock$ accepts inputs $\mathbf{x} \in \mathbb{R}^d$.
	It returns logits $K \in \mathbb{R}^{h \times k}$ given by
	$K = \left(\sum^{d}_{j=1}W_{\cdot,\cdot,j}x_j\right)
	     \odot \left(\sum^{d}_{j=1}V_{\cdot,\cdot,j}x_j\right)$.
	An optional noise term $\varepsilon_K$ may be added.
	Winner dropout sets $\mathrm{argmax}_k(K_{h,\cdot})$ to $-\infty$
	with probability $p_{\mathrm{drop}}$, ramped $0 \to 0.1$ over the
	annealing horizon.

	Soft labels $P$ are given by $\mathrm{softmax}_{k}(K/T)$
	where $T$ is a temperature hyperparameter. It is annealed
	geometrically from $1.0$ to $0.03$ over the first 50k steps, then
	held, giving progressively sharper classifications.
	The architecture provides methods for causal interventions on labels.
	From $P$ we collect per-sample entropy $H$,
	between-sample similarity $\mathrm{cossim}_b$,
	and batch mean entropy $\mathrm{KL}_m$.
	In practice $H$ is not needed, as including a
	$\mathrm{cossim}_b$ loss also minimizes $H$.
	$\mathrm{KL}_m$ serves to anchor the batch mean of each label vector
	$\mathbf{p}_{i,\cdot}$ to a uniform distribution.
	This results in a uniform $\mathbf{p}_{i,\cdot}$ code corresponding
	to a maximally generic feature vector.

	A $\dictblock$ accepts labels 
	$P \in \mathbb{R}^{h \times k}$
	and returns features $F \in \mathbb{R}^{h \times e}$ given by
	$F = \sum^{k}_{j=1}\lvert D_{\cdot,j,\cdot}\rvert \, P_{\cdot,j}$.
	Because $\lvert D_{\cdot,j,\cdot}\rvert$ is an absolute value
	and $P_{hj} \in [0,1]$, features are strictly additive.
	From $F$ we collect between-head similarity $\mathrm{cossim}_h$ on
	the clean features, and the sparsity penalty $\mathrm{L1}$ after the
	noise $\varepsilon_F$ is added.
	Both $\mathrm{L1}$ and $\mathrm{cossim}_h$ tend to decrease with only
	$\mathrm{MSE}$ loss, but incorporating them into the loss speeds up convergence.
	$\mathrm{cossim}_h$ is of particular interest because it measures
	the extent to which heads encode nonoverlapping geometries.

	We add noise $\varepsilon_F$ to $F$, then sum over the head dimension to obtain
	pooled features $F' \in \mathbb{R}^{e}$, which are added to $R$.
}
\label{fig:dictenc}
\end{figure}

\begin{figure}[tp]
\centering
\begin{tikzpicture}[
    font=\footnotesize,
    >={Stealth[round,length=1.8mm,width=1.6mm]},
    lab/.style={inner sep=1.5pt, font=\scriptsize},
  ]
  %% ---- head i --------------------------------------------------------
  \fill[black!8]   (0,0) -- (2.4,0) -- (1.2,2.0) -- cycle;
  \draw[semithick] (0,0) -- (2.4,0) -- (1.2,2.0) -- cycle;
  \fill (0,0)     circle (1.5pt) node[below,lab]{$\mathbf{a}_{i1}$};
  \fill (2.4,0)   circle (1.5pt) node[below,lab]{$\mathbf{a}_{i2}$};
  \fill (1.2,2.0) circle (1.5pt) node[above,lab]{$\mathbf{a}_{i3}$};
  \fill (1.05,0.72) circle (1.7pt) node[left,lab]{$\mathbf{f}_i$};

  %% ---- head i' ------------------------------------------------------
  \fill[black!8]   (5.0,0) -- (7.4,0) -- (6.2,2.0) -- cycle;
  \draw[semithick] (5.0,0) -- (7.4,0) -- (6.2,2.0) -- cycle;
  \fill (5.0,0)   circle (1.5pt) node[below,lab]{$\mathbf{a}_{i'1}$};
  \fill (7.4,0)   circle (1.5pt) node[below,lab]{$\mathbf{a}_{i'2}$};
  \fill (6.2,2.0) circle (1.5pt) node[above,lab]{$\mathbf{a}_{i'3}$};
  \fill (6.55,0.60) circle (1.7pt) node[right,lab]{$\mathbf{f}_{i'}$};

  \draw[<->, densely dashed, semithick, black!55] (1.35,0.71) -- (6.30,0.61);
  \node[lab] at (3.85,1.35)
    {$\omega$ penalizes heads $i$, $i'$ sharing coordinates};

  \node[lab, anchor=north, align=center] at (3.7,-0.9)
  {$\mathbf{f}'=\sum^h_{i=1} \mathbf{f}_i$,\qquad
  $\mathbf{f}_i=\sum^k_{j=1} p_{ij}\,\mathbf{a}_{ij}$,\qquad $\sum^k_{j=1} p_{ij}=1$};
\end{tikzpicture}


\caption[Per-head geometry]{\texttt{DictBlock} \textbf{head geometry} for $h=2$; $k=3$.
	The output space of a head $i$ forms a simplex with the rows of the matrix
	$D_{i,\cdot,\cdot}$ as the vertices.
	The row vectors $\mathbf{d}_{i,j,\cdot}$ form the spanning set of the
	feature space. Any $F_i$ must be a linear combination of the rows of
	$D_{i,\cdot,\cdot}$.
	As $F$ becomes sparser, these vectors become more approximately
	orthogonal, allowing them to be treated as a basis.
	The label code $P$ corresponds to soft assignments to mutually
	contrastive categories rather than independently firing features.
	Because elements of $P$, $p_{ij} \in [0,1]$,
	outputs must be within the convex hull of $D_{i,\cdot,\cdot}$.
	Because $\lvert D\rvert\geq0$ and $P\geq0$, $F\geq0$;
	therefore $\cos(F_i, F_j) \in [0,1]$, and driving
	$\mathrm{cossim}_h(F)$ to zero demands \emph{disjoint support}
	between heads, strictly stronger than orthogonality of signed
	vectors. This is the sense in which the architecture is built to
	seek an orthogonal decomposition.}
\label{fig:head-simplex}
\end{figure}
